\section{Tabularray's libraries}

\subsection{Drawing in your table with Tikz}

\begin{warningbox}{Minimal tabularray version}
	The tikz library is only available from tabularray version 2025A or more recent.
\end{warningbox}

\begin{Code}[show-output]
% \UseTblrLibrary{tikz}  % in the preamble

\highLight{\begin{tblrtikzbelow}}
	\path[fill=blue!15,draw=blue!80, ultra thick, rounded corners]
	(table.north west) rectangle (table.south east);
\highLight{\end{tblrtikzbelow}}

\highLight{\begin{tblrtikzabove}}
	\draw[red!70, thick]
		(2-2.north west) -- (2-3.south east)
		(2-2.south west) -- (2-3.north east);
	\draw[color=green,thick]
		(h1-|v1) -- (h1-|v2) -- (h2-|v2) -- (h2-|v3) -- (h3-|v3) -- (h3-|v4)
		-- (h4-|v4) -- (h4-|v5) -- (h2-|v5) -- (h2-|v6) -- (h1-|v6);
\highLight{\end{tblrtikzabove}}

\begin{tblr}{
	colspec={ccccc},
	hlines={4pt},
	vlines={3pt},
	hline{1,Z} = {0pt},
	vline{1,Z} = {0pt},
}
	1-1 & 1-2 & 1-3 & 1-4 & 1-5 \\
	2-1 & some text & other text & 2-4 & 2-5 \\
	3-1 & 3-2 & 3-3 & 3-4 & 3-5
\end{tblr}	
\end{Code}

The tikz library allows to draw on (or under) the table, with any valid tikz instruction, once the size of each cell has been computed\footnote{Note that you can insert tikz images in a table cell. But these images will increase the cell's width.}. The environment \snippet{tblrtikzabove} allows to draw on top of the table, and \snippet{tblrtikzbelow} under the table. These environments must be defined before the table they relate to.

Each cell is given a named tikz node: the cell at row i and column j is named \snippet{(i-j)}. Similarly, each border intersection has a name: the intersection of hline i with vline j is named \snippet{(hi-|vj)}. 

\subsection{Diagonal separations with diagbox}

\begin{Code}[show-output]
% \UseTblrLibrary{diagbox}  % in the preamble

\begin{tblr}{
	colspec={c|ccc},
}
	\highLight{\diagbox{Day}{Roommate}} & Martin & Emma & Julie\\ \hline
	Monday & Dishes & & Cooking\\
	Tuesday & Cooking & Dishes & \\
	Wednesday & & Cooking & Dishes \\
	Thursday & Dishes & & Cooking \\ 
	Friday & Cooking & Dishes & \\
	Saturday & & Cooking & Dishes \\
	Sunday & & & \\
\end{tblr}
\end{Code}

\snippet{\textbackslash diagboxthree\{\}\{\}\{\}} is also available.

%\subsection{Automatic in-cell computing with functional}
